% SAME-FAMILY SOURCE DRAFT, 25 October 2026.
% Author: audit_finite57_algorithms. Complete translation batch, not accepted.
% No imports, workers, enumeration, tests, Lean or TeX jobs.
% Root/different-author mathematical review and integration are pending.
\subsection{\(B_{17}\) with an actual isolated short spectator}
\label{app:B17I18}
We separate the local payment lemmas from the short-tail classification.
This makes the order of the argument explicit: the local lemmas first
dispose of the exceptional \(B_2\) cap, and only afterwards do we prove
the unconditional strong-five bootstrap. None of the local lemmas assumes
that bootstrap or an already strong terminal wall.

Throughout, \(P_i=e_1+\cdots+e_i\),
\(f_i=\sqrt i-.61/\sqrt2\), and \(\phi=(1+\sqrt5)/2\).
The original vector is \(z=E(p-x-ye)\), with
\(\norm x^2+y^2<1\). The actual isolated \(I\)-wall is
\[
 c_I=\sqrt h\,p-\sqrt{1+h}\,e,\qquad 0<h<1/2.
\]
Put
\[
 s^2={\sqrt{1+h}-1\over\sqrt{1+h}+1},\qquad
 t=\sqrt{1+4h}-2\sqrt h,\qquad h(r)={4r\over(1-r)^2}.
\]
Appendix~\ref{app:Ijoint} proves \(y\ge s\), and height at least
\(t\) in every other actual unit dual direction. At a classical long
end with \(\Omega_1=\sqrt2e_1\), this says \(x_1\ge t\).
In particular
\[
 0<s<1/3,\qquad (1-s^2)(1-t^2)=8st,\qquad h=h(s^2).
\]
These facts concern actual walls and the same original vector.
The local lemmas below require only their stated long-root, point and
reference hypotheses; a \(B\) short root is not among those hypotheses.
Where the initial roots are used, they are the actual long roots
\(n_j=(e_j-e_{j+1})/\sqrt2\), \(j=1,2,3\), with actual
long duals \(\Omega_1=\sqrt2P_1,\Omega_2=\sqrt2P_2\).

\begin{lemma}[The additional real \(I\)-barrier]
\label{lem:I923-real-barrier}
The actual parameters satisfy \(t+\tfrac52s>.923\).
\end{lemma}
\begin{proof}
The test value \(t_0=.923-2.5s\) is positive. Substitution in the
left-minus-right of the relation gives
\[
 P(s)=.148071-2.769s+13.601929s^2-4.615s^3+6.25s^4.
\]
For \(v=s-.1\in(-.1,7/30)\), its coefficients are
\(.00320029,-.1620642,12.592429,-2.115,6.25\).
Since \(-2.115v^3\ge-.4935v^2\),
\[
 P(s)>.0032-.163|v|+12v^2>0.
\]
The last discriminant is \(.026569-.1536<0\). The original
expression decreases strictly with positive \(t\), proving the claim.
This is separate from the earlier \(.767\)-barrier.
\end{proof}

\subsubsection{Local paid squares, with every initial cross retained}
Call \(W_0\) a paid classical vector if its demand \(D_0\) has already
been obtained from nonnegative actual-pressure coefficients. Its effective
lower energy is \(\mathcal E_0=2D_0-\norm{W_0}^2\).
At a retained anchor \(P_m\), let the actual demand be \(a_m\);
any debit below is made from its remaining individual coefficient.
No \(I\)-energy is included in \(\mathcal E_0\).

\begin{lemma}[A paid head square at an arbitrary initial anchor]
\label{lem:I-head-paid-square}
Suppose an actual pure classical pressure is
\(\ip xZ\ge H\), with \(Z\) supported in the first \(m\) coordinates
and \(\ip Z{P_m}=1\). Write \(R=Z^2\) and
\[
 B_Z=\ip{W_0}Z+{a_m-\ip{W_0}{P_m}\over m}.
\]
For \(\nu\ge0\), with debit \(\nu/m\) available at \(P_m\), the
same square improves \(\mathcal E_0\) by
\[
 2\nu(H-B_Z)-\nu^2(R-1/m).
\]
\end{lemma}
\begin{proof}
Use \(W_0+\nu(Z-P_m/m)\), after removing \(\nu/m\) from the
anchor coefficient. Its residual square is \(R-1/m\); expanding its
full cross with \(W_0\) gives the formula. In particular this formula
does not assume a weak-five \(P_1/P_6\) reference is flat.
\end{proof}

\begin{lemma}[A joint point and the \(I\)-height on a paid reference]
\label{lem:I-joint-paid-reference}
Suppose the actual pressures are \(a x_1+v y\ge F\), \(y\ge s_0\),
with \(a>0,v\ge0\). Put
\[
 B_0=(W_0)_1+{a_m-\ip{W_0}{P_m}\over m}.
\]
For \(\eta,\tau\ge0\), write \(k=\eta a\), \(\ell=\eta v+\tau\),
and require the individual anchor debit \(k/m\) to be available.
The vector
\(W_0+k(e_1-P_m/m)+\ell e\) improves \(\mathcal E_0\) by
\[
 2\eta(F-aB_0)+2\tau s_0-(1-1/m)k^2-\ell^2.
\]
If \(\ell=s_0\), put \(X=(F-vs_0)/a\). The improvement becomes
\[
 s_0^2+2k(X-B_0)-(1-1/m)k^2,
 \qquad \tau=s_0-kv/a\ge0.
\]
\end{lemma}
\begin{proof}
Remove \(k/m\) at the original anchor before using the two new
pressures. Expand the square, retaining
\(\ip{W_0}{e_1-P_m/m}\). The old vector has no \(e\)-component.
This proves both identities and keeps the nonnegative multiplier and
individual-price conditions visible, also when \(k\) is capped.
\end{proof}

\begin{lemma}[A capped initial graft at a weak-five reference]
\label{lem:I-capped-initial-graft}
Let \(a_0=.5965\). Suppose \(W_0\) is already paid, its first
coordinate is \(a_0\), and \(\ip{W_0}{P_6}\) equals its retained
\(P_6\) demand. If the remaining \(P_6\) price is at least \(1/60\),
then
\[
 \lambda=\min\{(6/5)(t-a_0)_+,.1\}
\]
gives, using \(e_1-P_6/6\) and the single unused \(I\)-height,
an improvement
\[
 s^2+2\lambda(t-a_0)-(5/6)\lambda^2>.015.
\]
\end{lemma}
\begin{proof}
The full old-reference cross changes the baseline to \(a_0\);
it is not zero in general. In the uncapped branch, weighted Cauchy
and Lemma~\ref{lem:I923-real-barrier} give
\[
 s^2+(6/5)(t-a_0)_+^2>
 {(.923-.5965)^2\over25/4+5/6}>.015.
\]
In the capped branch \(t>a_0+1/12\). The relation makes \(t\)
decrease with \(s\), and
\[
 t(.1)={\sqrt{11401}-40\over99}<67/99<a_0+1/12.
\]
Thus \(s<.1\). The \(.923\)-barrier then gives
\[
 s^2+.2(t-a_0)-1/120>
 (s-.25)^2-83/15000>.0169666.
\]
The debit is at most \(1/60\). These are one paid square and one
\(I\)-height, not two independent initial-coordinate gains.
\end{proof}

A flat-five reference below means a nonnegative actual-anchor reference
\(w\), with demand \(\norm w^2=q\), first five entries
\(b=f_5/5<.361\), and remaining \(P_5\) coefficient \(>.146\).
The anchor demand there is \(f_5=5b\). Extra terminal anchors are allowed.
Every local gain below preserves their nonnegative coefficients.
A head means the actual opposite of the old long node one.

\begin{lemma}[The foreign-cleared head scalar]
\label{lem:I-cleared-head-flat-gain}
On a flat-five reference, suppose an actual pure pressure has
\[
 Z=(1-u/2)e_1+(u/2)e_2,\quad
 H={\sqrt{2-u}-\sqrt{u(1-u)}\over\sqrt2},\quad
 R=1-u+u^2/2,\qquad 0<u<1.
\]
Weight \(\nu=.36\), costing \(.072\) at \(P_5\), gives gain \(>.04\).
\end{lemma}
\begin{proof}
Use Lemma~\ref{lem:I-head-paid-square}. Put
\(.5091<.36\sqrt2<.5092\). Lower chords for \(\sqrt{2-u}\)
and upper chords for \(R\), on the three successive intervals, are
\[
\begin{array}{c|c|c}
[0,.5]&1.4142-.379u&1-.75u\\
[.5,.75]&1.4381-.4268u&.8125-.375u\\
[.75,1]&1.4721-.4721u&.625-.125u.
\end{array}
\]
Their endpoint squares and concavity/convexity prove the bounds.
After \(b<.361\), the gains exceed
\[
\begin{array}{c|c}
.35636922-.0958u-.5092\sqrt{u(1-u)}&.04936922\\
.39283671-.16873u-.5092\sqrt{u(1-u)}&.04022171\\
.43444611-.2242u-.5092\sqrt{u(1-u)}&.04414611.
\end{array}
\]
The right column follows from
\(pu+c\sqrt{u(1-u)}\le(p+\sqrt{p^2+c^2})/2\),
using upper radical seals \(.5182,.5365,.5564\).
\end{proof}

\begin{lemma}[Finite heads on a flat-five reference]
\label{lem:I-finite-head-flat-reference}
Suppose the retained classical long-end diagram has only an isolated
head or a simple-three extension at old node two. On any flat-five
reference, a nonisolated head gives gain \(>.0994\).
For a pure finite isolated head the head gain alone is \(>.0186\),
and its gain including the single \(I\)-height is \(>.0933875\).
For a positive foreign finite head the combined gain is \(>.0957\).
Consequently \(q\ge.90715\) excludes every nonisolated or finite
isolated head beside actual \(I\), without a terminal-strength hypothesis.
\end{lemma}
\begin{proof}
For the nonisolated head,
\[
 A^2=2g^2+2g+Q,\qquad C_1={2g+1\over\sqrt{1-2Q}}.
\]
Own zero has zero time or the positive null
\(c+\Omega_2/2=\sqrt Qp-\beta e\);
the dual has its actual positive old-root expansion.
Own three has the triangle and parallel the pair null.
Every other own gives \(x_1>.643\). Replacing the first five entries
by \(.643,(5b-.643)/4,\ldots\) is legal: its debit
\((.643-b)/4\) fits the individual \(P_5\) coefficient, and its
gain exceeds \((5/4)(.282)^2>.0994\).

For an isolated finite head
\(c=Ap-g\sqrt2e_1-\beta e\), \(g<1\), put \(Q=\beta^2\).
The distinct actual \(I\)-wall gives
\(\lambda=\sqrt h A-\sqrt{1+h}\beta\ge0\).
At \(\beta=0\), \(0<\lambda<1/\sqrt2\), so \(\lambda=1/2\).
At \(\beta>0\),
\[
 \lambda<\sqrt{(1+\beta^2)/2}-\beta\sqrt{3/2}
 \le(\sqrt{10}-\sqrt6)/4<1/2.
\]
The middle upper bound decreases with \(\beta\ge1/2\), and its
bound before substitution increases with \(h<1/2\); hence \(\lambda=0\).
Squaring gives \(\beta^2=h(2g^2-1)<1/2\), so \(\beta=1/2\).
Both cases consequently have
\[
 h(2g^2-1)=1/4,\qquad h>1/4.
\]
Own zero/three/four is impossible, five gives \(h=\phi/2>1/2\),
six gives \(h=1/2\), and parallel has \(c+n_1\).
Thus a finite isolated own has label at least seven.

If \(\beta>0\), \(u=(1-Q)/g^2<1\). The actual positive core
\[
 c+\beta e+gu n_1
\]
has products zero on \(c,e\), product \(g(u-1)<0\) on \(n_1\),
square \(-g^2u(1-u)\), and positive time. It yields exactly the
pure pressure of Lemma~\ref{lem:I-cleared-head-flat-gain}.
The \(I\)-height costs \(>.0557\), by \(h>1/4\), so the gain
including \(I\) is \(>.0957\).

If \(\beta=0\), put \(u=g^{-2}\in(1,1.232)\). The actual core
\[
 c+g^{-1}n_1+2(g^{-1}-g)n_2
\]
has zero cap/first-root products, negative second/third-root products
and positive time. Its pure data are
\[
 Z=(1-u/2)(e_1+e_2)+(u-1)e_3,\quad
 H=\sqrt{1-u/2}-\sqrt{(7u-3u^2-4)/2},\quad
 R=3-4u+3u^2/2.
\]
The native-seven cubic gives \(g>.90096\), and
\(1.232(.90096)^2>1\). Weight \(1/2\) gives
\[
\begin{array}{c|c|c|c}
u&\text{head gain lower}&h\text{ lower}&I\text{ energy lower}\\ \hline
[1,1.1]&.0477&1/4&.0557\\
[1.1,1.15]&.033025&11/36&.066\\
[1.15,1.2]&.0213875&23/68&.072\\
[1.2,1.232]&.0186&3/8&.079 .
\end{array}
\]
The first radicals are bounded below by
\(.6708,.6519,.63245,.6196\), the second above by
\(.1871,.204125,.204125,.2\), and \(R\) above by
\(.5,.415,.38375,.36\). Positive squaring proves these;
the second radicand has maximum \(1/24\) and \(R\) decreases.
Use \(h=u/[4(2-u)]\) and \(h(r)\) for the \(I\)-column.
Every combined gain is \(>.0933875>1-.90715\).
Every head pressure has positive actual-root coefficients, and all
unselected original walls retain nonpositive products.
\end{proof}

\begin{lemma}[An actual joint point for every pure dotted mutual]
\label{lem:I-pure-dotted-joint-point}
For an actual pure dotted head \(c=Ap-g\sqrt2e_1\), put
\[
 u=g^{-2}\in(0,1),\quad h_0={u\over4(2-u)},\quad
 a=\sqrt{1-u/2},\quad v_0=\sqrt{u/(8-3u)},\quad
 d=\sqrt{1-a^2-v_0^2},\quad F=1-d,\quad s_0=s(h_0).
\]
For \(u\ge1/5\), every actual \(h\ge h_0\) gives a feasible
future point \((p-ae_1-v_0e)/d\). Consequently
\[
 a x_1+v_0y\ge F,\quad y\ge s_0,\qquad
 X={F-v_0s_0\over a}
 =\sqrt{{1-\sqrt{3u/8}\over1+\sqrt{3u/8}}}.
\]
\end{lemma}
\begin{proof}
The actual mutual product is \(-A\sqrt h<0\), hence
\(h\ge1/(4A^2)=h_0\). Put \(v=\sqrt{h/(1+h)}\);
then \(v\ge v_0\), \(v<1/\sqrt3\).
The denominator \(D=u/2-v_0v\) is positive because
\(3u(8-3u)\ge111/25>4\) on this interval. For
\[
 T=1/D,\quad \rho={Ta\over\sqrt2g}>0,\quad
 \sigma={Tv_0\over\sqrt{1+h}}>0
\]
direct substitution gives
\[
 p+\rho c+\sigma c_I=T(p-ae_1-v_0e).
\]
Its \(c\)-product is zero and its \(c_I\)-product is
\(T\sqrt{1+h}(v_0-v)\le0\). Old wall products are nonpositive
by duality; all other actual walls are nonpositive termwise.
Its time is positive and \(d^2>0\). Thus the normalized point is
actual even if the actual \(c,c_I\) Gram is not positive.
The identities \(v_0s_0=1-\sqrt{1-v_0^2}\) and
\(d/\sqrt{1-v_0^2}=\sqrt{3u/8}\) give \(X\).
It decreases with \(u\), whereas \(s_0^2\) increases.
The comparison \(X>r\) at \(u=U\) is the positive rational test
\[
 8(1-r^2)^2-3U(1+r^2)^2>0.
\]
\end{proof}

\begin{lemma}[The foreign dotted-head parabola]
\label{lem:I-foreign-head-parabola}
For a dotted head with \(Q=\beta^2>0\), set \(v=(1-Q)/g^2\).
Then \(0<v<3/4\), and
\[
 h\ge h_{\min}(v):={v\over6-4v}.
\]
The actual foreign-cleared core gives the pure pressure
\[
 Z=(1-v/2)e_1+(v/2)e_2,\quad
 H(v)=\sqrt{1-v/2}-\sqrt{v(1-v)/2},\quad
 R(v)=1-v+v^2/2.
\]
For \(F(v)=.8H(v)-.16R(v)\), \(\zeta=v-2/3\),
\[
 F(v)\ge C_0+A_1\zeta+.448\zeta^2,\quad
 C_0>.2976415,\quad0<A_1<.008385,\quad F(v)>.2976.
\]
Also \(y^2>.08625v-.014375\) for \(v\le.6\), and
\(y^2>.0385\) for \(v\ge.6\).
\end{lemma}
\begin{proof}
Acuteness with \(c_I\) gives \(h(2g^2-1)\ge Q\); hence
\[
 h\ge{Qv\over2(1-Q)-v}\ge{v\over6-4v}.
\]
The second expression increases with \(Q\ge1/4\).
The core \(c+\beta e+((1-Q)/g)n_1\) has zero cap/old-\(I\)
products, negative \(n_1\)-product, square \(-g^2v(1-v)\)
and positive time. It clears the foreign coordinate before producing
the displayed pure pressure; every other wall is nonpositive.

Put \(r=\sqrt6\) and \(\ell=r/3-r\zeta/8-.09\zeta^2\).
On this interval \(\ell\ge7r/24-.04>0\).
The exact identity
\[
 1-v/2-\ell^2
 =\zeta^2(3r/50-3/32-(9r/400)\zeta-.0081\zeta^2)
\]
has nonnegative right side, since the bracket is at least
\(.045r-.09735>0\). Further,
\((\sqrt{v(1-v)/2}-1/3)^2\ge0\) gives
\[
 H(v)\ge(r-1)/3+(2-r)\zeta/8+.66\zeta^2.
\]
Thus \(C_0=4(r-1)/15-4/45\) and
\(A_1=19/75-r/10\). Use
\(2.449489<r<2.449490\) and
\(.008385^2<1.792(.2976415-.2976)\).

For \(v\le.6\), \(h_{\min}\le1/6\), and
\(.23(2.081)^2<1\) gives \(s(h_{\min})^2>.23h_{\min}\).
The exact tangent inequality is
\[
 {v\over6-4v}-\left({3v\over8}-{1\over16}\right)
 ={3(2v-1)^2\over8(6-4v)}\ge0.
\]
For \(v\ge.6\), \(h_{\min}\ge1/6\), whose \(I\)-energy is
\(>.0385\), again checked using \(h(r)\).
\end{proof}

\begin{lemma}[The head on any sufficiently strong flat-five reference]
\label{lem:I-head-flat-reference}
Assume the actual long-end whitelist of
Lemma~\ref{lem:I-finite-head-flat-reference}.
On any flat-five reference every actual head gives a legal gain,
including \(I\) once, greater than \(.07665\).
In particular \(q\ge.9240775\) excludes all heads.
No actual strong-fifteen condition is required.
\end{lemma}
\begin{proof}
Finite and nonisolated heads are covered by
Lemma~\ref{lem:I-finite-head-flat-reference}. For a pure dotted head
with \(u\le.2\), its actual core \(c+g^{-1}n_1\) gives the
pressure \(Z,H,R\) of Lemma~\ref{lem:I-cleared-head-flat-gain}.
Here \(H>.6658,R\le1\). Weight \(.3\), with price \(.06\),
gives gain \(.6(.6658-.361)-.09(.8)=.11088\).

For pure \(.2\le u\le.5\), use weight \(\eta=.33\) of the
joint point alone. In Lemma~\ref{lem:I-pure-dotted-joint-point},
\[
 F>.5839,\quad a<.949,\quad (4/5)a^2+v_0^2<.748.
\]
The first bound follows at \(u=.5\), since \(d^2\) increases;
the last follows at \(u=.2\), since that expression decreases.
The gain is
\[
 .66(.5839-.949\cdot.361)-.1089(.748)=.07780806.
\]
Its individual price is \(<.066\); no \(I\)-energy is added again.

For \(.5\le u\le.8\), use
Lemma~\ref{lem:I-joint-paid-reference} with \(m=5,B=.361\),
\(k=5(X-B)/4\), \(\ell=s_0\). The unused \(I\) multiplier
is nonnegative because \(X-B<.269<.4a\) and \(s_0/v_0\ge1/2\).
Its price is \(<.06725\). Positive rational tests give
\[
\begin{array}{c|c|c|c}
u&s_0^2\text{ lower}&X\text{ lower}&\text{gain lower}\\ \hline
[.5,.6]&.0200&.597&.08962\\
[.6,.7]&.0254&.567&.078445\\
[.7,.75]&.0315&.553&.07758\\
[.75,.775]&.0349&.547&.078145\\
[.775,.8]&.0366&.540&.07665125.
\end{array}
\]
Use the right endpoints for \(X\), and left endpoints in \(h(r)\)
for \(s_0^2\); for the last two these have \(h_0=3/20,31/196\).
For \(.8\le u<1\), the actual pure core gives \(H>.491,R\le.52\).
Indeed \(H\) increases there and \(R\) decreases. A half-payment
gains \(>.05\), and \(h_0\ge1/6\) contributes \(>.0385\).

For a positive foreign dotted contact use
Lemma~\ref{lem:I-foreign-head-parabola} at weight \(.4\).
The head gain is \(>F-.2568\), at price \(.08\).
For \(v\le.6\), adding the \(I\)-height gives
\[
 >.0839665+(A_1+.08625)\zeta+.448\zeta^2>.0789665,
\]
using \((.09464)^2<1.792(.005)\).
For \(v\ge.6\) the gain is \(>.2976-.2568+.0385=.0793\).
All prices fit the single available \(P_5\) coefficient, and every
parameter branch gives more than the stated \(.07665\).
\end{proof}

\subsubsection{A conditional actual-helper packet at long index fifteen}
\begin{lemma}[Actual helper pressures and the retained fifteen packet]
\label{lem:I-helper15-pressure-packet}
Suppose an actual pure far-three opposite at a long index \(i\ge5\)
has
\[
 c_i=A_i p-\gamma\sqrt2P_i-\Omega_1/2,\qquad
 A_i^2=2i\gamma^2+2\gamma-1/2,
\]
where \(\sqrt2P_i\) is the actual old long dual,
\(n_i=(e_i-e_{i+1})/\sqrt2\), and the old initial roots
\(n_1,n_2\) are actual. Its own is assumed to be five, finite at
least seven, or dotted; the invoking local classification must establish
that list. Write \(H_k=P_k/(k\sqrt2)\).
The following are actual pure pressures \(\ip xU\ge T\):
\[
\begin{array}{c|c|c}
\text{own}&U&T\\ \hline
5&\gamma(\sqrt2P_i-n_i)+H_3&H(A_i,2/3-\gamma^2)\\
7,8&\gamma(\sqrt2P_i-n_i)+H_2&A_i-\sqrt{\gamma^2-3/4}\\
\text{finite }\ge9&\gamma(\sqrt2P_i-n_i)+H_1&H(A_i,1-\gamma^2)\\
\text{dotted}&\gamma\sqrt2P_i-\gamma^{-1}n_i+H_1&A_i-\sqrt{1-\gamma^{-2}},
\end{array}
\]
where \(H(A,d)=A/(1+\sqrt{d/(A^2+d)})\).

In particular suppose \(i=15\) is weak, and the legal reference \(w\)
ends at strong fourteen, is flat-five with \(q>.90715\), and has
individual \(P_{14}\) coefficient \(>.1361\).
Put \(\alpha=\sqrt2\gamma\) and
\[
 \Pi U=\alpha P_{14}+P_5/(5\sqrt2),\quad
 Y=U-\Pi U,\quad \Delta=T-\alpha f_{14}-b/\sqrt2.
\]
Then \(w.Y=0\), and the following complete table is valid:
\[
\begin{array}{c|c|c|c|c}
\text{own}&\Delta\text{ lower}&Y^2\text{ upper}&\mu&
 2\mu\Delta-\mu^2Y^2\text{ lower}\\ \hline
5&.39558&.724&.09&.06534\\
7&.340518&.963&.09&.05349294\\
8&.2873002&1.005&.09&.043573536\\
\text{finite }\ge9&.299602&1.4&.09&.04258836\\
\text{dotted}&.3307&1.68&.08&.04216 .
\end{array}
\]
Thus \(W_0=w+\mu Y\) has effective energy \(>q+.0421\).
Its debits are \(<.1273\) at \(P_{14}\) and \(<.01273\) at \(P_5\).
Its initial coordinate is
\[
 (W_0)_1=b+{\mu(1/k-1/5)\over\sqrt2}<.412.
\]
This lemma does not assume a \(B\) short normal or that \(P_{16}\)
is an actual long dual.
\end{lemma}
\begin{proof}
For own five clear old \(n_i,n_1,n_2\), with positive correction
\(\gamma n_i+2n_1/3+n_2/3\); the actual finite block is \(H_4\).
Its residual is \(2/3-\gamma^2>0\), so use the finite projected
point, not the corrected spacelike root. For owns seven/eight clear
\(\gamma n_i+n_1/2\); its residual is \(3/4-\gamma^2<0\).
For finite owns at least nine clear \(\gamma n_i\), and for dotted
owns use \(c_i+n_i/\gamma\). Selected products are zero or negative,
time is positive, and all other actual walls are nonpositive termwise.
These include the old \(I\), its opposite, any long head and either
classical tail. The displayed \(U\)'s follow by subtraction in the
original Euclidean coordinates.

At fifteen the initial residual is
\(H_k-P_5/(5\sqrt2)\); the remaining residual has support only in
fifteen/sixteen. It is orthogonal to every retained reference anchor.
For finite owns,
\[
 Y^2=\gamma^2+1/(2k)-1/10;
\]
for dotted owns it is \(2\gamma^2-1.6+\gamma^{-2}\).
Use \(\sqrt{28}-.61<4.682\), \(b/\sqrt2<.256\).
For finite five/seven/eight, the respective
\((A_i\text{ lower},\text{radical upper},\gamma\text{ upper})\)
are
\[
 (4.555,.111,.81),\quad(5.064,.249,.901),\quad
 (5.191,.322,.9239).
\]
Their positive squares give the first three deficits.
For finite \(\gamma\ge.939\), the deficit increases:
\(H_A>.99\), \(A_i'>\sqrt{30}\), \(H_d<0\), and
\(d=1-\gamma^2\) decreases. At \(.939\),
\[
 A_i>5.275,\quad\sqrt{A_i^2+d}>5.286,\quad\sqrt d<.344,
 \quad H>4.952,
\]
giving \(\Delta>.299602\).

For dotted fifteen, weakness implies \(\gamma<10/9\):
at \(10/9\) its squared edge ratio is
\(10609/4320>(39/25)^2\), whereas \(T(\sqrt{30})<39/25\).
The deficit decreases on this interval, since
\(A_i'<5.6\), the radical derivative is \(>1\), and
\(\sqrt{28}-.61>4.681\). At \(10/9\),
\(A_i^2=2093/54>6.225^2\) and the radical is \(<.436\);
thus \(\Delta>.3307777\). The norm is \(<1.68\).
All entries and prices now follow by rational multiplication.
The largest initial coordinate occurs for \(k=1,\mu=.09\),
and is below \(.361+.09(.8)/\sqrt2<.412\).
\end{proof}

\begin{lemma}[A retained helper, a dotted head and actual \(I\)]
\label{lem:I-helper15-head-coupling}
Assume the actual helper and reference hypotheses of
Lemma~\ref{lem:I-helper15-pressure-packet}, with \(f_{14}<3.31034\),
and an actual isolated-within-classical dotted head. For a pure head
with \(u\ge.2\), a single paid square gives helper gain \(>.0421\)
and joint head/\(I\) gain \(>.0515\). For a positive foreign head
the corresponding gains are \(>.0421\) and \(>.0521\).
For a pure head with \(u\le.2\), the standalone gain \(>.11088\)
on \(w\) suffices without paying the helper.
In particular each case contradicts the original norm budget.
\end{lemma}
\begin{proof}
For a pure head with \(u\ge.2\), take \(B=.412\) above the actual
initial coordinate \(B_0\). In
Lemma~\ref{lem:I-joint-paid-reference}, use \(m=5\),
\(k=5(X-B)/4\), \(\ell=s_0\). Since the helper residual has
zero sum on the first five entries, its full cross changes the
baseline to precisely \(B_0\); it is not orthogonal to a head.
The multipliers are nonnegative: on \([.2,.4]\),
\(X<.756,a>\sqrt{.8}\), and on \([.4,1]\),
\(X<.665,a>1/\sqrt2\). Thus \(X-B<.4a\), while
\(s_0/v_0\ge1/2\). The total \(P_5\) price is
\(<.086+.01273<.146\), and the separate \(P_{14}\) price fits.
The complete continuous table is
\[
\begin{array}{c|c|c|c}
u&s_0^2\text{ lower}&X\text{ lower}&
 s_0^2+(5/4)(X-.412)^2\text{ lower}\\ \hline
[.2,.5]&0&.629&.05886125\\
[.5,.6]&.0200&.597&.06278125\\
[.6,.7]&.0254&.567&.05543125\\
[.7,.8]&.0315&.540&.05198\\
[.8,.9]&.0385&.514&.051505\\
[.9,1]&.0464&.490&.054005 .
\end{array}
\]
The \(X\)-tests are those of
Lemma~\ref{lem:I-pure-dotted-joint-point}; the \(I\)-tests use
the left endpoints \(h_0=1/12,3/28,7/52,1/6,9/44\).
For example the first \(X\)-test has positive numerator
\(.0002777057265\). This proves gain \(>.0515\) on all the
displayed intervals. Small \(u\) uses the standalone actual core
already proved in Lemma~\ref{lem:I-head-flat-reference}.

For a foreign head put \(v=(1-Q)/g^2\), and use its cleared
pressure at weight \(\nu=.4\). The initial helper products with
\(Z-P_5/5\), for \(k=3,2,1\), are respectively
\[
 {2/15\over\sqrt2},\quad {3/10\over\sqrt2},\quad
 {.8-v/2\over\sqrt2}.
\]
The last is largest on \(0<v<.75\). Hence the complete adverse
cross is at most \((.072/\sqrt2)(.8-v/2)\).
Using \(1/\sqrt2<.7072\), the head gain is
\[
 >F(v)-.29753472+.0254592v.
\]
For \(v\le.6\), adding the single \(I\)-height and the parabola gives
\[
 >.06020458+(A_1+.1117092)\zeta+.448\zeta^2
 >.05215458,
\]
since \(.1201^2<1.792(.00805)\).
For \(v\ge.6\) the gain is \(>.0538408\).
The combined \(P_5\) price is \(<.09273\); the \(P_{14}\) price
is unchanged. No \(I\)-energy was included in the pure cleared
pressure or the helper gain. The full totals exceed
\[
 .90715+.0421+.0515=1.00075,\qquad
 .90715+.0421+.0521=1.00135 .
\]
\end{proof}

\subsubsection{The fresh \(B_{17}\) tail and its exceptional actual mate}
We now impose the entire \(B_{17}+I\) forest, with actual unit roots
\(n_i=(e_i-e_{i+1})/\sqrt2\), \(i\le16\), and genuine short
root \(n_{17}=e_{17}\). Thus
\(\Omega_i=\sqrt2P_i\) through sixteen and \(\Omega_{17}=P_{17}\).
All classical opposites retain the old \(I\).
The valid bulk parameter is \(M=12\), giving pure-helper coverage
through thirteen and the actual mixed-point nonadjacency.

At fourteen deletion retains \(A_{13}+B_3\).
Its only right finite extensions are long-three at fifteen and
short-three at seventeen, the latter the sealed \(F_4\).
Their \((\Delta,\sigma)\) are
\((1/2-Q,\gamma+1/2)\) and
\((1/4-Q,\gamma+1/(2\sqrt2))\).
Own zero has the actual affine star or \(3,4,3,3\) path;
own three on the long side is a triangle, and all other positive
cases are directly strong. Full bridges attach by three at left
endpoints and the long right endpoint, are finite \(B_{17}\) and
foreign-sealed. A short-end full bridge would put the old four
internally in a long path. Left/isolated/foreign cases retain the
complete long-root inventory of Appendix~\ref{app:prefixinventory}.
Consequently coverage and nonadjacency reach fourteen, before
assuming that five is strong.

At fifteen the retained right block is the actual \(B_2\) on
long sixteen and short seventeen. A finite new contact is three
at either end only. The data are
\[
 (\Delta,\sigma)=(1/2-Q,\gamma+1/2),\quad
 (1/2-Q,\gamma+1/(2\sqrt2)).
\]
The long case has its own-zero star and own-three triangle.
In the short case, own zero has the affine \(F_4\) kernel
\((\sqrt2,2\sqrt2,3,2,1)\) on \(c,n_{17},n_{16},n_{15},n_{14}\).
Own four has \(\sqrt2c+n_{14}+2n_{15}+n_{16}\),
own six has \(\sqrt3c+2n_{15}+n_{14}\), and parallel has
\(c+n_{15}\). For long own zero use
\(c+n_{15}+2n_{16}+\sqrt2n_{17}\).
Direct products give zero selected products, positive cap time
and nonpositive products with every remaining original wall.
Every other own at least five is strong. Positive foreign
\(Q\ge1/4\) makes the short own-three case strong too.
All four left/right endpoint bridges are retained, with residual
\(1/30\); their minimum ratio, including own zero, is
\[
 \sqrt{30}\bigl(1/30+1/(2\sqrt2)\bigr)>100/61.
\]
The full bridges are foreign-sealed. The sole additional weak cap is
\[
 c=Ap-(\Omega_{15}+\Omega_{17})/2,\qquad Q=0.
\]

Release old sixteen and keep this same actual cap. The retained
path old one through fifteen, then \(c\), then the actual short
seventeen is finite \(A_{17}\), all bonds three. Positive finite
projection is future and satisfies every original wall. Omitted
sixteen is strict, because it has negative product with the positive
old-fifteen and short-seventeen terms. This is the actual other
endpoint of old edge sixteen. Its exact data give
\[
 C_{16}={16+15\sqrt2\over\sqrt{18}},\qquad
 S_{16}\ge4\sqrt{{C_{16}-1\over C_{16}+1}}
 =G:={8+12\sqrt2\over7},\quad
 3.56722<G<3.5673<f_{16}.
\]
In particular this floor is not called strong sixteen.

\begin{figure}[ht]
\centering
\setlength{\unitlength}{1mm}
\begin{picture}(106,34)
\put(4,25){\makebox(0,0){\(\cdots\)}}
\put(12,25){\circle*{1.4}}
\put(12,25){\line(1,0){26}}
\put(38,25){\circle*{1.4}}
{\color{gray}\put(38,25){\line(1,0){26}}
\put(64,25){\circle{2.4}}
\put(64,25){\line(1,0){26}}
\put(77,29){\makebox(0,0){\(4\)}}
}
\put(90,25){\circle*{1.4}}
\put(38,25){\line(2,-1){26}}
\put(64,12){\circle*{1.4}}
\put(64,12){\line(2,1){26}}
\put(12,31){\makebox(0,0){\(n_{14}\)}}
\put(38,31){\makebox(0,0){\(n_{15}\)}}
\put(64,31){\makebox(0,0){\(n_{16}\)}}
\put(90,31){\makebox(0,0){\(n_{17}\)}}
\put(64,6){\makebox(0,0){\(c\)}}
\end{picture}
\caption{The exceptional cap replaces old sixteen. The solid path
through \(n_{15},c,n_{17}\) has only simple bonds. The grey old edges
are released; the old four bond is not changed into an invented root.}
\label{fig:B17Iactualmate}
\end{figure}

If five is weak, the actual first-six pressure gives
\(x_1>a_0=.5965,S_6>2.01815\).
The reference through one, six and the actual floor \(G\) at sixteen
has decreasing slopes and energy
\[
 >.7600299945+(3.5672-2.01815)^2/10
 =.99998558475.
\]
Its \(P_6\) coefficient exceeds \(.28433-.154915>1/60\).
Lemma~\ref{lem:I-capped-initial-graft} therefore gives a contradiction.
This uses no head reduction or later strong-five conclusion.

If five is strong, take the already available strong thirteen/fourteen
reference and append the straight chord to \(G\). Its slopes
\(<.132,<.129\) fit the individual terminal coefficients.
The new energy is respectively
\[
 >.88864+.39299^2/3>.9401,\qquad
 >.90715+.25688^2/2>.9401.
\]
For the energy bounds use \(f_{13}<3.17423,f_{14}<3.31034\).
For the slope test instead use \(f_{13}>3.1741,f_{14}>3.3102\)
and \(G<3.5673\), all given by positive squaring.
The first five block and its full \(P_5\) coefficient are unchanged.
Deleting the long head leaves actual \(B_{16}\), whose finite
one-node extension is only isolation or simple-three at long two.
Thus Lemma~\ref{lem:I-head-flat-reference} applies to this stronger
reference, independently of any strong-fifteen assertion.
This contradicts the original norm and excludes the exceptional cap
in both initial branches.

\subsubsection{Unconditional strong five from two retained actual helpers}
Coverage through fifteen now consists only of nonadjacent pure helpers.
If five is weak with finite own, the actual first-six table gives
\(x_1>.65\). With strong six and an actual strong fourteen/fifteen,
positive interpolation gives
\[
 .65^2+(2.01815-.65)^2/5+1.2921^2/8
 =1.00555718575>1.
\]
Thus a weak-five own would be dotted, \(1<a<6/5\).
If fifteen is strong, the first-six/strong-fifteen interpolation has
energy \(>.985\); its \(P_6\) capacity exceeds \(1/60\).
The capped initial graft adds \(>.015\), again impossible.
It remains that both five and fifteen are weak, so fourteen is strong.

Let \(w\) interpolate through \((1,a_0),(6,f_6)\) and every
actual strong knot through fourteen. Put \(b_1=(f_6-a_0)/5\);
then \(.28433<b_1<.284332\). The rational increment floors give
\[
 \norm w^2>.9703299945,\quad
 b_{P_1}>.312168,\quad b_{P_6}>.08806,\quad b_{P_{14}}>.1361.
\]
For example \(b_{P_6}>b_1-(\sqrt7-\sqrt6)\), and there are
at most four nonadjacent omissions among seven through thirteen.

Retain the complete actual dotted-five foot
\[
 U_5=\sqrt2aP_5-a^{-1}n_5+\Omega_1/2,\quad
 T_5=\sqrt{10a^2+2a-1/2}-\sqrt{1-a^{-2}}.
\]
Its projection onto the initial reference span is \(mP_6+lP_1\), where
\[
 m=4\sqrt2a/5<1.35766,\qquad
 l=\sqrt2a/5+1/\sqrt2<1.04652.
\]
The residual \(Y_5\) is orthogonal to \(w,P_1,P_6\), and
\[
 Y_5^2=8a^2/5-2+a^{-2}<1,\qquad T_5-\ip w{U_5}>.1203.
\]
Indeed \(\ip w{P_5}=a_0+4b_1<1.733828\). The deficit
decreases with \(a\): the time derivative is \(<13/4\),
the radical derivative is \(>1\), and
\(\sqrt2\ip w{P_5}>2.4519\).
At \(a=6/5\), use
\(A>4.03732,\sqrt{11}/6<.552771,\sqrt2<1.414214\).
Weight \(\mu_5=.05\) consequently gives gain \(>.00953\),
with debits \(<.067883\) at \(P_6\) and \(<.052326\) at \(P_1\).

For the complete weak-fifteen packet, project its initial \(H_k\)
onto \(P_1,P_6\). The projection of \(U_{15}\) is
\[
 \alpha P_{14}+l_kP_1+m_kP_6,\qquad
\begin{array}{c|cc}
k&l_k&m_k\\ \hline
3&\sqrt2/10&\sqrt2/15\\
2&\sqrt2/5&\sqrt2/20\\
1&1/\sqrt2&0 .
\end{array}
\]
The residual \(Y_{15}\) is orthogonal to all original reference
anchors. Here \(\ip w{H_3}<.275,\ip w{H_2}<.312,
\ip w{H_1}<.422\). The proof of
Lemma~\ref{lem:I-helper15-pressure-packet}, with these initial
demands in place of \(b/\sqrt2\), gives
\[
\begin{array}{c|c|c|c|c}
\text{own fifteen}&T_{15}-\ip w{U_{15}}\text{ lower}&
Y_{15}^2\text{ upper}&\mu_{15}&\text{gain lower}\\ \hline
5&.37658&.824&.09&.06111\\
7,8&.2313&1.105&.09&.0326835\\
\text{finite }\ge9&.1336&1&.09&.015948\\
\text{dotted}&.1647&1.28&.08&.01816 .
\end{array}
\]
All points remain actual; the table uses the pure-helper fifteen
classification established before this bootstrap.

The full mixed crosses are
\[
 \ip{Y_5}{Y_{15}}=
 \begin{cases}2a/15&k=3,\\a/10&k=2,\\0&k=1.\end{cases}
\]
Their adverse square terms are at most \(.00144,.00108,0\).
The first two rows already exceed \(1.0395,1.0114\), without
any \(I\)-energy. Their combined \(P_6\) debit is
\(<.067883+.008486<.08806\).

For the last two rows, both residuals are orthogonal to
\(e_1-P_6/6\). The already paid vector still has first coordinate
\(a_0\) and \(P_6\) demand \(f_6\), so apply
Lemma~\ref{lem:I-capped-initial-graft}.
The total \(P_6\) debit is
\(<.067883+1/60<.08455<.08806\); the \(P_{14}\) debit
is \(<.1273\), and \(P_1\) debits are \(<.117<.312168\).
The full original-reference/graft cross is retained in that lemma.
One square, with \(I\) once, therefore exceeds
\[
 .9703299945+.00953+.015948+.015=1.0108079945,
\]
or
\[
 .9703299945+.00953+.01816+.015=1.0130199945.
\]
Every weak-five case is impossible. Five is now unconditionally strong.

\subsubsection{The final composition}
Take a strong fourteen/fifteen reference. It has \(q>.90715\),
first-five mean \(b\), and individual \(P_5\) coefficient \(>.146\).
The actual \(B_{16}\) long-head whitelist, and
Lemma~\ref{lem:I-finite-head-flat-reference}, exclude all finite
or nonisolated heads. The remaining actual head is isolated-within-\(B\)
and dotted, retaining every finite contact with old \(I\).

If fifteen is strong, its reference has \(q>.9240775\), so
Lemma~\ref{lem:I-head-flat-reference} excludes every such head.
If fifteen is weak, fourteen is strong and its actual cap is a pure
helper, so Lemma~\ref{lem:I-helper15-head-coupling} applies.
Its small-parameter branch uses only the standalone actual head point;
its other branches retain the full initial helper/head cross and
the single \(I\)-height. They give the stated strict norm contradictions.
This excludes the entire \(B_{17}+\)actual-\(I\) row.

All local points were positive combinations or positive finite
projections of original walls, with every selected sign and future
orientation checked before use. The original farthest vector is never
replaced. No arbitrary face is declared Coxeter, no disconnected root
scale or arithmeticity is imposed, no actual mate floor is promoted to
strong sixteen, and no earlier elimination energy is added again.

